\chapter{Точная LCF-проверка связывающего GKSL полугруппы}
\label{ch:v8-linking-qms-gksl-lcf-migration-gate}

\section{Типизированный генератор}

На точном концевом пространстве
$E=\mathbb C^{11}\oplus\mathbb C^{10}$ физическая матрица инцидентности $A$
имеет только целые элементы $0,1$ и ранг $10$. Самосопряжённый оператор
\begin{equation}
 D_A=\begin{pmatrix}0&A^*\\A&0\end{pmatrix}
 \label{eq:v8-lcf-linking-dirac}
\end{equation}
является корректным эндоморфизмом $E$. Конструктор GKSL принимает нулевой
гамильтониан, один оператор скачка $D_A$ и точную скорость $1$, получая
\begin{equation}
 \mathcal L(Z)=D_AZD_A-\frac12\{D_A^2,Z\}
 =-\frac12[D_A,[D_A,Z]].
 \label{eq:v8-lcf-linking-gksl}
\end{equation}

Сохранение следа проверено точно на полном базисе из $21^2=441$ матричной
единицы; $\mathcal L(I_{21})=0$ проверено как точное матричное равенство.

\section{Алгебра концевых секторов и явная формула}

Для всех $121+100=221$ матричных единиц блочно-диагональной алгебры
$M_{11}\oplus M_{10}$ образ снова блочно-диагонален. На том же полном базисе
LCF-ядро проверяет равенство
\begin{equation}
 \mathcal L(X,Y)=\left(
 A^*YA-\frac12\{A^*A,X\},
 AXA^*-\frac12\{AA^*,Y\}
 \right).
 \label{eq:v8-lcf-linking-corner-formula}
\end{equation}
Следовательно, формула исходного гейта подтверждена как тождество линейных
отображений, а не как выборочная численная проверка.

\section{Точная неподвижная размерность}

Неподвижность блочно-диагонального наблюдаемого эквивалентна системе
\begin{equation}
 AX=X_tA,
 \qquad
 XA^*=A^*X_t.
\end{equation}
Её точная матрица имеет форму $220\times221$, ранг $180$ и нульмерность
\begin{equation}
 221-180=41.
 \label{eq:v8-lcf-linking-fixed-forty-one}
\end{equation}

\begin{theorem}[Точный связывающего GKSL сертификат]
Физическая матрица инцидентности задаёт типизированный конечномерный GKSL-
генератор, точно сохраняющий след и единицу, оставляющий
$M_{11}\oplus M_{10}$ инвариантной алгеброй и имеющий на ней неподвижное
пространство размерности $41$.
\end{theorem}

\begin{nogocriterion}[Граница полной положительности]
Полная положительность полугруппы в текущем eDSL следует из доверенного
конечномерного правила GKSL. Её нельзя называть независимо формализованной,
пока ядро не проверяет отдельный Чой/Краус объект доказательства для $e^{t\mathcal L}$
или само правило не перенесено в Lean.
\end{nogocriterion}

\begin{protocol}
\begin{enumerate}
 \item размерность носителя концевых секторов: \textbf{$21$};
 \item ранг инцидентности: \textbf{$10$};
 \item GKSL-типизация: \textbf{пройдена};
 \item матричных единиц проверки следа: \textbf{$441$};
 \item унитальность: \textbf{точно};
 \item базис концевой алгебры: \textbf{$221$};
 \item явная угловая формула: \textbf{точно};
 \item коммутантная система: \textbf{$220\times221$};
 \item ранг и нульмерность: \textbf{$180$ и $41$};
 \item независимый Чой-сертификат: \textbf{ещё нет};
 \item следующий перенос: \textbf{усреднение по калибровочной группе Краус-мост}.
\end{enumerate}
\end{protocol}

Машинный сертификат сохранён в
\path{s2t_v8_linking_qms_gksl_lcf_migration_gate_results.json}.